In this section we prove \Cref{theorem:certificate iff non-empty
language}. The proof is split into the two directions. In
\Cref{ssec:certificate-to-reachability} below, we show that the existence of a witnessing
certificate for an \laVASS implies that the language of the \laVASS is
non-empty. The converse direction is then shown
in \Cref{ssec:non-empty-to-certificate}.

\subsection{Witnessing certificates imply language non-emptiness}\label{ssec:certificate-to-reachability}

This section proves the following proposition.
%
\begin{proposition}\label{lem:abstract-to-concrete}
  If there exists a witnessing certificate for a restricted \laVASS
  $V$ then $L(V)\neq \emptyset$.
\end{proposition}
%

The idea behind the proof of \Cref{lem:abstract-to-concrete} is that
we obtain from a witnessing certificate $(R,X,Y,L)$ of an \laVASS $V$
a sequence of runs of $V$ such that the final run in that sequence is
an accepting run of $V$. Initially, we obtain a run that ends in a
configuration where the counters are in a congruence relation. We then
carefully pump the simple loops pointed to by $L$, beginning from the
last counter working towards the first.

To formally prove \Cref{lem:abstract-to-concrete}, let $(R,X,Y,L)$ be
a witnessing certificate, and let $\pi(R)$ be the run in the
configuration graph of $V$ induced by $R$. Let $a\le d$ be maximal
such that $\ell_a\neq \top$. We now define a sequence of runs
$\pi_0,\ldots,\pi_a$ such that the following invariant holds. In the final
configuration of $\pi_i$,
%
\begin{enumerate}[(i)]
\item $m_j \le n_j$ and $m_j\equiv n_j \bmod \ell_j$ for the $j$-th
  counter pair, $1\le j \le a-i$; and
\item $m_j = n_j$ for the $j$-th counter pair, $a-i<j\le d$.
\end{enumerate}
%
It is clear that $\pi_a$ then witnesses $L(V)\neq \emptyset$. We
proceed by induction on $i$.

\emph{Base case $i=0$:} Let $\pi_0=\pi(R)$. Since $R$
is a witnessing certificate, $val(\pi(R),x_a)\leq val(\pi(R),y_a)$, and hence
$m_a \le n_a$ in the last configuration of $\pi_0$.
%Since $R$ is a
%witnessing certificate, we also have $m_{a+1}=n_{a+1}, \ldots,
%m_d=n_d$ by definition.
Moreover, $R$ respects the set of induced $y$-constraints. Hence
$n_{a-1} - n_a \ge \delta_{a-1}$, where $\delta_{a-1}$ is the length
of the path from $R[X(a-1)]$ to $R[X(a)]$. Hence $n_{a-1} - n_a \ge
m_{a-1} - m_a$ and thus $m_{a-1} \le n_{a-1}$. Iterating this argument
for the remaining counters, we get that~(i) of the invariant is
fulfilled for $\pi_0$; (ii) trivially holds since $R$ ends in an
accepting abstract configuration.
%

\emph{Induction step  $i>0$:} Let 
$\pi_{i-1}$ be the path that exists by the induction hypothesis. If
$m_{a-i}=n_{a-i}$ in the last configuration of $\pi_{i-1}$ then we are
done and take $\pi_i=\pi_{i-1}$; otherwise $m_{a-i} < n_{a-i}$ and
$m_{a-i} \equiv n_{a-i} \bmod
\ell_{a-i}$. Hence, there is some $k\in \N$ such that $n_{a-i}= k
\cdot \ell_{a-i}$. Since $\ell_i\neq\top$, let
$\beta\defeq\alpha_{L(a-i)} \xrightarrow{t_1} \alpha_2
\xrightarrow{t_2} \cdots
\alpha_{\ell_i-1}\xrightarrow{t_{\ell_i}}\alpha_{L(a-i)}$ be the
simple $\alpha$-loop at position $L(a-i)$ that is guaranteed to exist
since $R$ is a witnessing certificate. We insert the transitions of
$\beta^k$ and the induced updated configurations into $\pi_{i-1}$ at
position $L(a-i)$. Notice that $L(a-i) < X(a-i+1)$. Otherwise, by the
definition of the induced $y$-constraints,
$y_{a-i}-y_{a-i+1}=\delta_{a-i}$ is in the set of induced
$y$-constraints, where $\delta_i= X(a-i+1)-X(a-i)$. Since the last
abstract configuration of $R$ respects the set of $y$-constraints, it
must be the case that in the last configuration of $\pi_{i-1}$,
$n_{a-i}-n_{a-i+1}= \delta_{a-i}$ and
$m_{a-i}-m_{a-i+1}=\delta_{a-i}$, so $m_{a-i}=n_{a-i}$, because after
the position $X(a-i+1)-1$ in $R$ and thus $\pi_{i-1}$, the counters
$x_{a-i},x_{a-i+1}$ get incremented simultaneously. This contradicts
our assumption that $m_{a-i}\neq n_{a-i}$. Thus, the counters
$x_{a-i+1},y_{a-i+1},\ldots,x_d,y_d$ remain unchanged by the insertion
of $\beta^k$, so (ii) and consequently~(i) continues to hold in the
last configuration of $\pi_i$ for those counters. Moreover, due to the
ordering conditions imposed on witnessing certificates, the value of
$y_{a-i}$ does not change either, and hence $m_{a-i}=n_{a-i}$ in the
last configuration of $\pi_i$. Since $\beta$ is a loop in the abstract
configuration space, we have $m_j \equiv n_j \bmod \ell_j$ for all
$1\le j< a-i$ and the values of $u_j$, for all $1\le j<a$ are
preserved.


\subsection{Reachability yields witnessing certificates}\label{ssec:non-empty-to-certificate}

We now turn towards the converse direction and show that we can obtain
a witnessing certificate from a run witnessing non-emptiness. 
%
\begin{proposition}\label{lem:concrete-to-abstract}
  If a restricted \laVASS $V$ admits an accepting run then
  there exists a witnessing certificate for $V$.
\end{proposition}
%
We begin with defining a function that turns a configuration from
$C(V)$ into an abstract configuration. This function is parameterised
by $\ell_1,\ldots,\ell_d\in \N_+ \cup \{ \top \}$:
%
\begin{multline*}
f_V((q,m_1,n_1,\dots,m_d,n_d),\ell_1,\dots,\ell_d)\defeq 
(q,m_1\circ \ell_1,n_1\circ \ell_1,\dots,m_d\circ \ell_d,n_d\circ \ell_d,\\
min(n_1-n_2,U_1),\dots,min(n_{d-1}-n_{d},U_{d-1}),\ell_1,\dots,\ell_d)\,.
\end{multline*}
%
Here, $m\circ \ell\defeq\bot$ if $m=0$; $m\circ \ell\defeq
m\bmod \ell$ if $\ell\in\N_+$; and $m\circ \ell\defeq m$ if
$\ell=\top$.  We lift the definition of $f_V$ to paths of concrete
runs $\pi$ in the natural way, and write
$f_V(\pi,\ell_1,\dots,\ell_d)$ for the resulting sequence of abstract
configurations. Let $\pi=c_1 \xrightarrow{t_1}
c_2 \cdots \xrightarrow{t_{n-1}} c_n$ be a run witnessing
$L(V)\neq\emptyset$. We show how to obtain a witnessing certificate
$R$ from $\pi$. Without loss of generality, in $c_n$ we have $m_1\ge
m_2 \ge
\ldots m_d>0$.

To this end, we show how from the accepting run $\pi$ we can
iteratively define a sequence $R_0,R_1,R_2\ldots,R_d$ of abstract runs
and identify the required $\ell_1,\ldots, \ell_d \in
\N_+ \cup \{ \top \}$ and $X,Y,L$ such that $(R_d,X,Y,L)$ is a reachability
certificate. Let $X(i)\defeq j$ such that $j$ is the first position in
$\pi$ where the value of counter $x_i$ is non-zero; analogously define
$Y(i)$ to be the first position where the value of $y_i$ is
non-zero. Clearly, $X,Y$ are monotonic and $X(i)\le Y(i)$, for all
$1\le i\le d$. Otherwise, if a counter $y_i$ gets initialised before
the counter $x_i$ in $\pi$, it must be the case that $n_i>m_i$ in
$c_n$ and therefore $c_n$ cannot be an accepting configuration.

Recall that $\pi$ has length $n$. In our proof, the subsequent
technical lemma will allow us to conclude that, if for a counter pair
$x_i, y_i$ the $y_i$ counter gets updated shortly after the $x_i$
counter then the run will end shortly after and counter pairs
$x_j,y_j$ for $j\ge i$ will consequently have small values.
%
\begin{lemma}\label{lm:smallpath}
If $Y(i)-X(i)\le dM_i$ for some $1\le i\le d$ then $n-Y(i)< dM_i$, so
$m_j,n_j\le 2dM_i$ in $c_n$ for all $i\le j\le d$.
\end{lemma}
\begin{proof}
We have that $Y(i)-X(i)\le dM_i$ implies that $val(\pi[1,Y(i)],x_i)\le dM_i+1$, and since 
$val(\pi[1,Y(i)+k],y_i)\ge 2^k$ we get that:
\begin{itemize}
\item $val(\pi,y_i)\ge 2^{n-Y(i)};$ and
\item $val(\pi,x_i)\le dM_i+n-Y(i)+1.$
\end{itemize}
Assume that $n-Y(i)\ge dM_i$. Then, $2^{n-Y(i)} - (dM_i+n-Y(i)+1)> 2^{n-Y(i)} - (2n-2Y(i)+1) >0$, if $n-Y(i)\ge 3$. However, $\pi$ is an accepting path, so $val(\pi,x_i)=val(\pi,y_i)$, and we get a contradiction. Thus, we must have that $n-Y(i)< dM_i$ which implies that $val(\pi,x_i)\le 2dM_i$, so $m_i=n_i\le 2dM_i$ and for any 
$j$, $i< j\le d$, $m_j\le m_i$ and $n_j\le n_i$, so $m_j,n_j< 2M_i$ 
in $c_n$, for all $i\le j\le d$ since $\pi$ is an accepting path.
\end{proof}

Let $R_0 \defeq f_V(\pi,1,1,\ldots,1)$. Note that $R_0$ together with
$X$ and $Y$ as defined above adheres to Conditions~(a)--(e) of
reachability certificates.
%% It remains to either identify that
%% $\ell_1=\top$ or to identify a suitable loop in $R_0$. We distinguish
%% two cases:
%% %
%% \begin{itemize}
%% \item $Y(1)-X(1) \le d\cdot M_1$: we choose $\ell_1 \defeq \top$ and
%%   $L(1)\defeq X(1)$.
%%   %% \Cref{lm:smallpath} guarantees that
%%   %% $(f_V(\pi, \ell_1,\ldots,\ell_d),X,Y,L)$ is a witnessing
%%   %% certificate.
%% \item $Y(1)-X(1) > d\cdot M_1$: then
%%   there is a segment in $R_0[X(1),Y(1)]$ of length greater than $M_1$
%%   on which no $x$-counter has its first $++$ transition. Hence, there is a
%%   smallest $k$, $X(1)\le k<Y(1)$ and a $\alpha_k \xrightarrow{t_{k}}
%%   \cdots \xrightarrow{t_{k+\ell}} \alpha_{k+\ell+1}=\alpha_k$ in $R_0$
%%   since $M_1>\abs Q$. In particular $k+\ell+1\le Y(1)$. We choose
%%   $L(1)\defeq k$ and $\ell_1 \defeq \ell$.
%% \end{itemize}
%% %
%% We define $R_1 \defeq f_V(\pi,\ell_1,1,\ldots,1)$.

Suppose $R_{i-1}$ and $\ell_1,\ldots,\ell_{i-1}$ have been
constructed. If $i>1$, and $L(i-1)\ge X(i)$ or $\ell_{i-1} = \top$
then we choose $\ell_{i} \defeq \ell_{i-1}$, $L(i)=L(i-1)$ and
$R_i \defeq f_V(\pi,\ell_1,\ldots,\ell_{i},1,\ldots,1)$. Otherwise, we
distinguish two cases.
%
\begin{itemize}
\item $Y(i)-X(i) < dM_i$: we choose $\ell_{i} \defeq \top$ and $L(i)\defeq X(i)$.
\item $Y(i)-X(i) \ge dM_i$: then there is a segment in $R_{i-1}[X(i),Y(i)]$ of length greater than $M_i$ on which no $x$-counter has its first $\wPP$ transition. Let $N_i$ be the maximum number of different abstract configurations on this segment. Since $\ell_{i-1}\neq\top$ we know that $m_j,n_j$ can take at most $M_j$ different values for all $1\le j<i$, as they can either be $\bot$ or a residue class modulo $M_j$. Also, for all $i\le j\le d$ the values of $m_j,n_j$ have a constant value, either 0 or $\bot$,
on this segment, and $u_i=\cdots=u_d=0$ in all abstract configurations
of this segment. So
%
\begin{align*}
  N_i \le{} & \abs Q \prod_{1\le j<i} M_j^2 \cdot U_j\\
  \le{} & \abs Q \prod_{1\le j<i} \abs Q^{(1/4)\cdot 32^{j-1} + 2 + 32^{j-1} + 4 }\\
  \le{} & \abs Q^{(1/3968)(5\cdot 32^i - 23968)+6i+1}\\
  <{} & \abs Q^{(1/8)\cdot 32^{i-1} + 1}\\
  ={} & M_i 
\end{align*}
%
By the pigeonhole principle, there is a smallest $k$, $X(i)\le k<
Y(i)$, $\ell < M_i $, and a simple loop
$\alpha_k \xrightarrow{t_k}\cdots \xrightarrow{t_{k+\ell}}\alpha_{k+\ell+1}=\alpha_k$
in $R_{i-1}$. We choose $L(i)\defeq k$, $\ell_i\defeq \ell$ and let
$R_i\defeq f_V(\pi,\ell_1,\ldots,\ell_{i},1,\ldots,1)$.
\end{itemize}
%
By construction, $(R_d,X,Y,L)$ is a reachability certificate. It
remains to turn it into a witnessing certificate. In particular, this
requires to removes from $R_d$, to ensure that the final segment of
$R_d$ is short, and to establish that $R_d$ is consistent with the
implied $y$-constraints.

Let $R \defeq R_d = f_V(\pi,\ell_1,\ldots,\ell_{d})$ and $a$ be the
largest index such that $\ell_a\neq \top$. In order to make $R$
loop-free, we iterate the following process:
%
\begin{itemize}
\item identify the first simple loop $\alpha_k \xrightarrow{t_k} \cdots \xrightarrow{t_{k+\ell}}
  \alpha_{k+\ell+1}$ in $R[1,Y(a)]$ and replace it by $\alpha_k$;
  observe that for $I \defeq \{k+1,\ldots,k+\ell\}$, we have $I \cap
  \{ X(i), Y(i), L(i) : 1\le i \le d \} = \emptyset$ since
  $\alpha_{X(i)-1} \xrightarrow{t_{X(i)-1}} \alpha_{X(i)}$ occurring in $R_d$ means
  that $x_i$ has value $\bot$ in $\alpha_{X(i)-1}$ and a value different from $\bot$ in
  $\alpha_{X(i)}$, and thus $\alpha_{X(i)}$ cannot be part of a
  loop; the same argument applies to any $Y(i)$. Finally, since $L(i)$
  was chosen as the index of the first configuration of the first
  cycle appearing after $X(i)$, we have $L(i)\not\in I$ for all $1\le
  i\le d$ as well.
\item update $X,Y,L$ such that for all $i$ such that $X(i)>k$,
  $X(i)\defeq X(i)-\ell$, and analogously $Y(i) \defeq Y(i) - \ell$
  and $L(i) \defeq L(i) - \ell$ for the respective $i$.
\end{itemize}
%
This process guarantees that $R[1,Y(a)]$ is loop-free. It is easy
to verify that $(R,X,Y,L)$ obtained in this way is a reachability
certificate and that the last abstract configuration of $R$ is
accepting.

We now show that the $y$-constraints induced by $R$ are valid in the
final configuration of $R$. To this end, we first show that for all
$1\le i\le d$ such that $\ell_i\neq\top$, $X(i+1) - X(i) <
U_i$. Consider the simple path $\alpha_{X(i)}
\xrightarrow{t_{X(i)}} \alpha_{X(i)+1} \xrightarrow{t_{X(i)+1}} \cdots \xrightarrow{t_{X(i+1)-1}}
\alpha_{X(i+1)}$. If $Y(i) \ge X(i+1)$ then clearly $X(i+1)-X(i) \le
N_i<U_i$, where $N_i$ is defined as above. Otherwise, there is a
$k\in\N$ such that the path decomposes as
\[
\alpha_{X(i)} \xrightarrow{t_{X(i)}}\cdots \alpha_{Y(i)}\xrightarrow{t_{Y(i)}}\cdots 
\alpha_{Y(i)+k}\xrightarrow{t_{Y(i)+k}}\cdots \xrightarrow{t_{X(i+1)-1}}\alpha_{X(i+1)}
\]
and
\begin{itemize}
\item $u_i=0$ in all abstract states $\alpha_j$ with $X(i)\le j\le
  Y(i)$;
\item $u_i=U_i$ in all abstract states $\alpha_j$ with $Y(i)+k\le j\le
  X(i+1)$; and
\item $k\le \log{U_i}$.
\end{itemize}

Thus, the maximum length of $R[X(i),X(i+1)]$ is bounded by:
%
\begin{align*}
  & N_i \cdot M_i + \log U_i + N_i \cdot 2M_i\\
  \le {} & 2 \cdot M_i^3 + \log U_i\\
%  \le {} & \abs Q^{(6/7)\cdot 32^{i-1}+4} + \abs Q^{5(i-1)+1}+4\abs Q\\
  \le {} & \abs Q^{(3/8)\cdot 32^{i-1}+4} + \abs Q^{5(i-1)+1}+4\abs Q\\
  < {} & \abs Q^{32^{i-1}+4}\\
  = {} & U_i
\end{align*}

We can now show that $R$ respects the induced $y$-constraints. Fix
some $1\le i\le d$ such that $\ell_i\neq \top$. We distinguish two
cases:
\begin{itemize}
\item  There is no $1\le j\le a$ such that $X(i)\le L(j)\le X(i+1)$. Thus, we know that $y_i-y_{i+1}=\delta_i$ is in the set of induced $y$-constraints. Also, $val(\pi,y_i)-val(\pi,y_{i+1}) = val(\pi,x_i)-val(\pi,x_{i+1}) = X(i+1)-X(i) = \delta_i$ since we did not remove any abstract loops on the segment of $R_d$ between the first $\wPP$ update for $x_i$ and the first $\wPP$ update for $x_{i+1}$. Finally, since $\delta_i<U_i$ by the above argument, we conclude that $u_i=\delta_i$ in the last abstract
configuration $R[n]$ of $R$. 
\item Otherwise, $X(i)\le L(i)\le Y(i)$, so $y_i-y_{i+1}\ge\delta_i$ 
is in the set of induced $y$-constraints. However, $val(\pi,y_i)-val(\pi,y_{i+1}) = val(\pi,x_i)-val(\pi,x_{i+1}) \ge X(i+1)-X(i) = \delta_i$ and again because $\delta_i<U_i$ we can conclude that $u_i\ge \delta_i$ in $R[n]$.
\end{itemize}
%
This establishes that the $y$-constraints are satisfied.  Let $n$ be
the index of the last abstract configuration of $R$. For the final
step, we now argue that $val(\pi(R),x_a)\le val(\pi(R),y_a)$ and $n-Y(a)\le
2dM_{d+1}$. We make a case distinction:
%
\begin{itemize}
\item $a=d$: Note that $val(\pi,x_d)=val(\pi,y_d)$.  Since we only
  remove loops from $R_d[1,Y(d)]$, we have that $val(\pi(R),x_d)\le
  val(\pi(R),y_d)$. If $n-Y(d)\le 2dM_{d+1}$ we are done with $(R,X,Y,L)$ as
  a witnessing certificate.  Otherwise, assume $n-Y(d)>
  2dM_{d+1}$. This implies that the path $R[Y(d),n]$ must contain at
  least one simple loop. Consider iterating the following process:
\begin{itemize}
\item remove the first simple loop from $R[Y(d),n]$ and update
  $n\defeq n-\ell$, where $\ell$ is the length of the simple loop that
  was removed; and                 
\item stop if $n-Y(d)\le 2dM_{d+1}$.
\end{itemize}
We argue that, $n-Y(d)\ge M_d^2$. Let $R'$ and $n'$ 
be the previous values of $R,n$ before the last iteration. It must be that $n'
>2dM_{d+1}$ and since the length of 
any simple loop of $R'[Y(d)+1,n']$ is bounded by 
$M_{d+1}$, we get that $n-Y(a)\ge M_{d+1}\ge M_d^2$. 
Note that $Y(d)-X(d)\le M_d\cdot \abs{Q}\cdot\prod_{1\le j< d}M_j^2\cdot U_j\le M_d^2$, 
so $val(\pi(R[1,Y(d)]),x_d)\le M_d^2$. It must be then 
the case that $val(\pi(R),x_d)\le val(\pi(R),y_d)$.

\item $a<d$: we know $n-Y(a)\le 2dM_{d+1}$ by \Cref{lm:smallpath}. Moreover, we must have that $val(\pi(R),x_a)\le val(\pi(R),y_a)$ since $val(\pi,x_a)=val(\pi,y_a)$ and we do not remove loops after the counter $y_a$ is incremented. 

\end{itemize}
